% !TEX program = tectonic
\documentclass[9pt,a4paper]{extarticle}
\usepackage[top=0.55cm,bottom=0.5cm,left=0.9cm,right=0.9cm]{geometry}
\usepackage[T1]{fontenc}\usepackage[utf8]{inputenc}\usepackage{lmodern}
\usepackage[french]{babel}
\usepackage{amsmath,amssymb}\usepackage{textcomp}\usepackage{xcolor}\usepackage{booktabs}\usepackage{array}
\newcolumntype{R}[1]{>{\raggedright\arraybackslash}p{#1}}
\usepackage{enumitem}\usepackage{multicol}\usepackage{graphicx}
\usepackage[hidelinks]{hyperref}\usepackage{url}\urlstyle{sf}
\graphicspath{{figs_nb11/}}
\definecolor{bleu}{RGB}{20,77,144}\definecolor{vert}{RGB}{20,110,60}
\definecolor{rouge}{RGB}{160,30,30}\definecolor{grisf}{RGB}{110,110,110}
\definecolor{or}{RGB}{150,105,0}
\newcommand{\fig}[2]{\begin{center}\includegraphics[width=\linewidth]{#1}\end{center}\vspace{-2.5mm}{\scriptsize\color{grisf} #2}\par\vspace{0.6mm}}
\pagestyle{empty}\setlength{\parindent}{0pt}\setlength{\parskip}{0.85pt}
\setlength{\abovedisplayskip}{2.5pt}\setlength{\belowdisplayskip}{2.5pt}
\setlength{\abovedisplayshortskip}{1pt}\setlength{\belowdisplayshortskip}{1pt}
\renewcommand{\arraystretch}{0.90}
\setlength{\columnsep}{6mm}\setlength{\columnseprule}{0.3pt}
\renewcommand{\columnseprulecolor}{\color{grisf!40}}
% une méthode = un bloc : numéro, titre, et l'idée en une ligne
\newcommand{\meth}[3]{\vspace{2.2mm}{\color{bleu}\hrule height 0.9pt}\vspace{0.7mm}%
  {\color{bleu}\bfseries #1~---~#2}\par{\footnotesize\color{grisf}\emph{l'idée :} #3}\par\vspace{0.9mm}}
\newcommand{\concret}{\par\vspace{0.9mm}{\color{rouge}\bfseries\footnotesize CONCRÈTEMENT}\par\vspace{0.2mm}}
\newcommand{\lecon}[1]{\vspace{0.4mm}{\scriptsize\color{vert}$\blacktriangleright$~\textbf{ce que ça apprend :} #1}\par\vspace{0.2mm}}
\newcommand{\obs}[1]{\vspace{0.2mm}\hangindent=3.4mm\hangafter=0{\color{bleu}\footnotesize\textbullet}~{\small #1}\par\vspace{0.2mm}}
\newenvironment{pts}{\begin{itemize}[leftmargin=3.4mm,itemsep=0.2mm,topsep=0.4mm,parsep=0pt,label={\color{bleu}\footnotesize\textbullet}]}{\end{itemize}}
\begin{document}

\begin{center}
{\large\bfseries Construire un portefeuille à partir des déclarations du Congrès}\\[0.5mm]
{\small Quatre méthodes, deux lentilles de mesure --- et ce qui reste quand on a tout contrôlé}
\end{center}
\vspace{0.4mm}\hrule\vspace{0.8mm}
{\footnotesize
\textbf{En bref.} On part de la méthode la plus naïve et on l'améliore pas à pas, en mesurant à chaque étape.
Le gain vient d'un seul geste --- \textbf{pondérer aux dollars déclarés}. Il subsiste, après contrôle du marché,
de la taille, de la valeur et du momentum, un $\alpha$ de \textbf{$+1{,}7$ à $+2{,}3$~\%/an} --- \textbf{$+1{,}5$
à $+2{,}1$ net de frais} --- avec des $t$ de $1{,}8$ à $2{,}0$. \textbf{Un $\alpha$ positif, qui résiste aux quatre
facteurs et au passage d'ordre, et qui reste à confirmer hors échantillon.}}
\vspace{0.8mm}\hrule\vspace{1.2mm}

\begin{multicols}{2}

% ═══════════ PROTOCOLE ═══════════
{\color{bleu}\bfseries 0 $\cdot$ Le protocole commun --- et ce qui ne l'est pas}\par\vspace{0.3mm}
\begin{pts}
\item \textbf{Données} : 134\,464 déclarations brutes $\to$ \textbf{113\,369 exploitées} (350 élus,
2\,755 titres, 2014--2026). Chaque déclaration donne un titre, un sens, une date, et une \emph{fourchette} de
montant --- on retient le \textbf{milieu} $A_k$.
\item \textbf{On se place à la divulgation} $\delta_k$, jamais à la transaction : c'est la seule date où
l'information est publique, donc la seule où la stratégie est réalisable. \emph{À la date de transaction tous les
chiffres seraient meilleurs --- et faux.}
\item \textbf{Recomposition mensuelle}, exécution au \textbf{close J+1}, \textbf{buy-and-hold} entre deux coupes :
voilà le socle commun. \emph{Ce qui n'est pas commun}, et qu'il faut avoir en tête : la \textbf{fenêtre} (42 jours
pour M1, 3 ans pour M2--M4) et le \textbf{périmètre} (M1 prend les deux chambres ensemble, sans distinguer les
partis ; M2--M4 traitent chaque parti \textbf{séparément et identiquement}).
\item \textbf{Trois mesures}, de la plus indulgente à la plus sévère :
\end{pts}
\[ \underbrace{x_Y=\!\prod_{t\in Y}\!(1{+}r_t)-\!\prod_{t\in Y}\!(1{+}r^{m}_t)}_{\text{excès brut}}
   \qquad t=\frac{\bar x}{\sigma_x/\sqrt n} \]
\[ \underbrace{r-r_f=\alpha_1+\beta(r^{m}\!-\!r_f)+\varepsilon}_{\text{on retire le marché}} \]
\[ \underbrace{r-r_f=\alpha_4+\beta_{\mathrm M}(r^{m}\!-\!r_f)+\beta_{\mathrm S}\mathrm{SMB}
   +\beta_{\mathrm H}\mathrm{HML}+\beta_{\mathrm U}\mathrm{MOM}+\varepsilon}_{\text{on retire aussi taille, valeur, momentum}} \]
\concret
{\small Les quatre facteurs ne sont pas de nous : ce sont ceux de Fama-French et Carhart, repris tels quels de la
bibliothèque de Kenneth French (\path{cache/ff_factors.csv}). Les trois derniers sont des \textbf{écarts entre
deux paniers} d'actions américaines --- on achète l'un, on vend l'autre :\par}
\begin{pts}
\item $r^{m}\!-\!r_f$ --- \textbf{le marché} : toutes les actions cotées, pondérées par capitalisation, moins le
taux sans risque.
\item \textbf{SMB} \emph{(Small Minus Big)} --- le pari sur la \textbf{taille} : acheter les petites entreprises,
vendre les grandes. Son rendement dit de combien les petites ont battu les grandes.
\item \textbf{HML} \emph{(High Minus Low)} --- le pari sur la \textbf{valeur} : acheter les actions bon marché
rapportées à leur valeur comptable, vendre les chères. Il dit de combien la valeur a battu la croissance.
\item \textbf{MOM} \emph{(Momentum)} --- le pari sur la \textbf{tendance} : acheter ce qui est monté sur l'année
écoulée, vendre ce qui a baissé. Il dit de combien les gagnants récents ont battu les perdants.
\end{pts}
\obs{\textbf{Comment lire un chargement.} Ces trois facteurs étant des \emph{écarts}, un chargement négatif se lit
à l'envers : $\beta_{\mathrm S}<0$ signifie que le portefeuille se comporte comme les \textbf{grandes}
capitalisations ; $\beta_{\mathrm H}<0$, comme les valeurs de \textbf{croissance}. Le $\alpha_4$ est ce qui reste
\textbf{une fois ces quatre expositions payées} --- c'est pour cela qu'il est le test le plus sévère du papier.}

% ═══════════ M1 ═══════════
\meth{M1}{Une voix par membre}{chaque élu qui achète a le même poids, quel que soit le montant --- l'hypothèse la
plus neutre possible.}
\concret
\begin{pts}
\item \textbf{Chaque fin de mois} $t$, on prend les opérations \textbf{divulguées} dans les 42 derniers jours de
bourse : $\ \mathcal K(t)=\{k:\ t-42<\delta_k\le t\}$.
\item Pour chaque titre, \textbf{les achats moins les ventes}, sans descendre sous zéro :
\[ B_i=\!\!\sum_{k\in\mathcal K,\ \text{achat}}\!\! A_k,\quad
   S_i=\!\!\sum_{k\in\mathcal K,\ \text{vente}}\!\!\gamma_k A_k,\quad
   \mathrm{Net}_i=\bigl[B_i-S_i\bigr]^{+} \]
\item \textbf{Mais une vente de titres détenus depuis plus de 2 ans ne compte pas.} On reconstitue les lots en
\emph{premier entré, premier sorti}, et $\gamma_k$ est la part de la vente portant sur des lots \textbf{récents} :
\[ \gamma_k=\frac{\sum_{h\le\bar a}q_h+r_k}{\sum_h q_h+r_k}\in[0,1],\quad \bar a=504\ \text{j}\approx2\ \text{ans} \]
$q_h$ = quantité vendue provenant d'un lot détenu depuis $h$ jours, $r_k$ = part \textbf{non appariée} (aucun achat
antérieur connu), comptée comme récente. Une vente qui solde une vieille position a $\gamma_k=0$ : elle ne dit rien
sur le titre aujourd'hui. \emph{En pratique} : $\gamma=1$ pour 80~\% des ventes, $\gamma=0$ pour 17~\%.
\item On ne garde que le solde positif, acheté par \textbf{au moins 2 élus} :
$\ \mathcal S_t=\{i:\ \mathrm{Net}_i>0,\ m_i\ge2\}$.
\item Chaque élu retenu a \textbf{une voix}, partagée également entre ses titres $\mathcal B_m$ --- s'il en a acheté
quatre, chacun reçoit un quart. \textbf{Le montant n'entre pas} : 5\,000~\$ pèsent autant que 500\,000~\$.
\[ w_i=\frac1{K_t}\sum_{m}\frac{\mathbf 1\{i\in\mathcal B_m\}}{|\mathcal B_m|},\qquad K_t=\#\{\text{élus acheteurs}\} \]
\item On achète au \textbf{close du lendemain} : $\ q_i=w_i\,V(t)/P_i(t{+}1)$, et on ne touche plus à rien
jusqu'à la coupe suivante --- les poids \textbf{dérivent} avec les cours.
\end{pts}
\begin{center}\scriptsize
\begin{tabular}{@{}l r r r r@{}}
\toprule
 & excès/an & $t$ & $\alpha_1$ & $\alpha_4$ \\
\midrule
toutes chambres & $+0{,}35$~\% & 0,18 & $-0{,}03$ & $+1{,}00$ \\
\bottomrule
\end{tabular}
\end{center}

% ═══════════ M2 ═══════════
\meth{M2}{Pondérer aux dollars déclarés}{celui qui engage 500\,000~\$ y croit plus que celui qui en engage 5\,000.}
\concret
\begin{pts}
\item Même rendez-vous mensuel, mais la fenêtre passe à \textbf{3 ans} et \textbf{on ne garde que les achats} ---
\textbf{les ventes sortent complètement du signal}, il n'y a plus de $\gamma$ ni de solde :
\[ \mathcal K^P(t)=\{k:\ P(k)=P,\ \text{achat},\ t-756<\delta_k\le t\} \]
\item Pour chaque titre, la somme des \textbf{dollars achetés} --- un achat de 500\,000~\$ pèse cent fois un achat
de 5\,000~\$ :
$\ B_i=\sum_{k\in\mathcal K^P,\ i(k)=i}A_k$.
\item Le poids est \textbf{proportionnel} à cette somme, sous \textbf{deux garde-fous} qui valent pour M2, M3
et M4 et ne changeront plus : on ne retient que \textbf{150 titres} et \textbf{aucune ligne ne dépasse 8,5~\%}.
Leur mécanique est détaillée au bloc suivant, où elle sert l'argument.
\[ w_i\ \propto\ B_i,\qquad\text{plafonné à }c=8{,}5\,\% \]
\item \emph{Pourquoi les poser dès maintenant} : la sélection des 150 se fait sur le score complet de M3. En la
gardant \textbf{identique} d'un bloc à l'autre, la comparaison M2 \emph{vs} M3 isole le seul facteur ajouté, et
jamais un effet de sélection. Ce qui n'est pas encore là, c'est le filtre de M4.
\end{pts}
\obs{\emph{Et si on nettait les ventes, comme le fait le fonds ?} Testé \textbf{sur M3} : l'excès NANC
\textbf{tombe} de $+2{,}98$ à $+1{,}40$~\%/an. Les écarter n'est pas une paresse, c'est ce qui marche.}
\begin{center}\scriptsize
\begin{tabular}{@{}l r r r r@{}}
\toprule
 & excès/an & $t$ & $\alpha_1$ & $\alpha_4$ \\
\midrule
NANC & $+1{,}35$~\% & 1,00 & $+0{,}31$ & $+0{,}96$ \\
GOP & $+0{,}52$~\% & 0,27 & $+0{,}21$ & $+1{,}28$ \\
\emph{repère : équipondéré} & \emph{$-0{,}67$~\%} & \emph{$-0{,}58$} & \emph{$-0{,}91$} & \emph{$+0{,}38$} \\
\bottomrule
\end{tabular}
\end{center}

% ═══════════ M3 ═══════════
\meth{M3}{Ajouter le consensus, et borner la concentration}{dix acheteurs valent mieux qu'un --- et aucune ligne
ne doit pouvoir emporter le portefeuille.}
\concret
\begin{pts}
\item On \textbf{multiplie} les dollars par le \textbf{nombre d'élus distincts} $m_i$ ayant acheté le titre.
\emph{Le geste se lit à somme égale} : 200\,000~\$ déclarés par \textbf{deux} élus donnent un score
\textbf{double} de 200\,000~\$ déclarés par un \textbf{seul}. On préfère l'argent dispersé à l'argent concentré.
\[ s_i=B_i\,m_i,\qquad m_i=\#\{m(k):k\in\mathcal K^P,\ i(k)=i\} \]
\item \emph{Conséquence, qu'il faut voir} : si $K$ élus déclarent chacun le même montant $a$, alors $B_i=Ka$
et $s_i=K^2a$ --- le score croît comme le \textbf{carré} du nombre d'élus, car $m_i$ est déjà dans $B_i$
(plus d'acheteurs, plus de dollars) et on le remultiplie.
\item Les \textbf{150 titres} annoncés en M2, ce sont les 150 plus gros \emph{scores} --- $\mathcal U(t)$ ---
et on normalise : $\ v_i=s_i\big/\sum_{j\in\mathcal U}s_j$.
\item On \textbf{plafonne chaque ligne à $c=8{,}5$~\%} : ce qui dépasse est redistribué au prorata sur les autres,
et on recommence jusqu'à ce que plus rien ne dépasse. Le point d'arrivée s'écrit d'un coup :
\[ \boxed{\,w_i=\min(c,\ \lambda v_i)\,}\qquad \lambda\ \text{tel que}\ \sum_i w_i=1 \]
$\lambda\mapsto\sum_i\min(c,\lambda v_i)$ croît de 1 à $|\mathcal U|c$ : $\lambda$ existe et est \textbf{unique}
dès que $|\mathcal U|\ge\lceil 1/c\rceil=12$ titres.
\item \textbf{Le seul changement par rapport à M2 est le facteur $m_i$} --- univers, plafond, fenêtre de 3 ans,
achats seuls, close du lendemain : tout le reste est identique. C'est ce qui rend la comparaison lisible.
\end{pts}
\obs{\textbf{La frontière du top-150 n'est pas nette} : les montants étant des \emph{fourchettes},
\textbf{104 titres} partagent le score du 150\textsuperscript{e} à une coupe de 2014 et 36 s'échangent selon
l'ordre de tri. \emph{Mesuré} : les trier de A à Z ou de Z à A déplace l'excès de \textbf{0,03 pt/an} --- ces
titres sont au score minimal, donc à poids quasi nul.}
\begin{center}\scriptsize
\begin{tabular}{@{}l r r c r@{}}
\toprule
 & excès/an & $t$ & IC 95~\% & $\alpha_4$ \\
\midrule
\textbf{NANC} & $+2{,}98$~\% & \textbf{1,47} & $[-1{,}4;+7{,}4]$ & $+1{,}75$ \\
\textbf{GOP} & $+1{,}24$~\% & \textbf{1,62} & $[-0{,}4;+2{,}9]$ & $+1{,}68$ \\
\midrule
\emph{$m_i$ seuls, sans les \$} & \emph{$+0{,}33$~\%} & \emph{0,33} & --- & \emph{---} \\
\bottomrule
\end{tabular}
\end{center}

% ═══════════ M4 ═══════════
\meth{M4}{Écarter le bruit des comptes gérés}{le gérant du fonds dit ignorer les dépôts en masse : \og{}\emph{300
trades \ldots{} 5 shares of something \ldots{} that's just a rebalance}\fg{}. On teste sa règle.}
\concret
\begin{pts}
\item \textbf{Avant tout calcul}, on compte combien d'opérations chaque élu a déclarées \textbf{le même jour} :
$\ n_k=\#\{k':\ m(k')=m(k),\ \delta_{k'}=\delta_k\}$.
\item On écarte celles qui appartiennent à un dépôt de \textbf{50 opérations ou plus} --- signature d'un compte
piloté par un courtier, pas d'une décision : $\ \mathcal K^P(t)$ est restreint à $\ n_k<\theta=50$.
\item Cela retire \textbf{54~\%} du \$ démocrate et \textbf{58~\%} du républicain.
\item Tout le reste est identique à M3.
\end{pts}
\begin{center}\scriptsize
\begin{tabular}{@{}l r r c r@{}}
\toprule
 & excès/an & $t$ & IC 95~\% & $\alpha_4$ \\
\midrule
NANC & $+4{,}55$~\% & 1,17 & $[-3{,}9;+13{,}0]$ & $+1{,}75$ \\
GOP & $+1{,}52$~\% & 1,04 & $[-1{,}7;+4{,}7]$ & \textbf{$+2{,}28$} \\
\bottomrule
\end{tabular}
\end{center}

\fig{meth_cascade.png}{la cascade, méthode par méthode : excès brut (gris), $\alpha$ marché seul (bleu),
$\alpha$ quatre facteurs (vert). Le passage du gris au bleu mesure ce que le marché explique.}

% ═══════════ LE TEST ═══════════
\meth{$\star$}{Le test qui tranche}{l'excès n'est-il qu'une exposition déguisée ? On régresse sur les quatre
facteurs du \S0.}
\begin{center}\scriptsize
\begin{tabular}{@{}l r r r r@{}}
\toprule
 & $\alpha_4$ & $t$ & $\beta_{\mathrm{SMB}}$ & $\beta_{\mathrm{HML}}$ \\
\midrule
M3 --- NANC & $+1{,}75$ & 1,81 & $-0{,}12$ & $-0{,}13$ \\
M3 --- GOP & $+1{,}68$ & 1,82 & $-0{,}08$ & $+0{,}06$ \\
M4 --- GOP & $+2{,}28$ & \textbf{1,99} & $-0{,}03$ & $+0{,}11$ \\
\emph{le SPY lui-même} & \emph{$+0{,}04$} & \emph{0,12} & \emph{$-0{,}12$} & \emph{$+0{,}01$} \\
\bottomrule
\end{tabular}
\end{center}
\fig{meth_facteurs.png}{les chargements factoriels. La ligne pointillée est le SMB du SPY ($-0{,}12$) : tous les
chargements de taille tiennent dans $|\mathrm{SMB}|\le0{,}12$.}
\lecon{\textbf{l'$\alpha$ survit au contrôle} --- il passe de $+1{,}26$ à $+1{,}75$ (NANC) au lieu de s'effondrer.
Le pari de taille qu'on redoutait \textbf{n'existe pas} : notre SMB \emph{est} celui du SPY. Et le SPY, passé dans
la même machine, rend $+0{,}04$ ($t$ 0,12) --- elle ne fabrique pas d'$\alpha$ à partir de rien. \textbf{Reste
que $t\approx1{,}8$} : positif, pas démontré.}

% ═══════════ LA STABILITÉ ═══════════
\meth{$\circlearrowright$}{La même mesure, année après année}{un $\alpha$ moyen sur douze ans peut n'être que deux
bonnes années. On rejoue la régression sur une fenêtre glissante d'un an --- et cette fois ce n'est pas le niveau
des courbes qui renseigne, c'est leur \emph{comportement}.}
\begin{center}\scriptsize
\begin{tabular}{@{}l r c r r@{}}
\toprule
 & $\beta$ méd. & $\beta$ min--max & $\alpha$ étendue & $\alpha>0$ \\
\midrule
M1 \emph{(commune)} & 1,00 & 0,90--1,08 & 29 pts & 57~\% \\
M2 --- NANC & 1,04 & 0,97--1,18 & 17 pts & 61~\% \\
M2 --- GOP & 0,97 & 0,83--1,22 & 24 pts & \textbf{51~\%} \\
\textbf{M3 --- NANC} & 1,05 & 0,96--1,26 & 23 pts & \textbf{70~\%} \\
\textbf{M3 --- GOP} & 0,99 & 0,89--1,13 & \textbf{14 pts} & \textbf{71~\%} \\
M4 --- NANC & 1,12 & 0,97--\textbf{1,49} & \textbf{48 pts} & 69~\% \\
M4 --- GOP & 0,97 & 0,84--1,12 & 17 pts & 69~\% \\
\bottomrule
\end{tabular}
\end{center}
\fig{meth_rolling.png}{$\beta$ (en haut) et $\alpha$ (en bas) réestimés sur chaque fenêtre d'un an ---
fenêtres chevauchantes : ces courbes \emph{décrivent}, elles ne testent pas.}
\lecon{\textbf{M3 est la seule régulière des deux côtés} : son $\alpha$ tient au-dessus de zéro 70~\% du temps
côté démocrate \emph{et} républicain, quand M2 y tombe à \textbf{51~\%} --- pile ou face. M4, lui, paie sa
performance en instabilité : $\beta$ de $0{,}97$ à $1{,}49$ d'un côté, $0{,}84$ de l'autre. \textbf{Et rien de
tout cela n'a servi à choisir M3} --- c'est une confirmation, pas l'argument.}

% ═══════════ LE FROTTEMENT ═══════════
\meth{\textdollar}{Ce que ça coûte à exécuter, et faut-il ralentir}{tous les chiffres précédents sont bruts.
Deux questions : combien coûte le passage d'ordre --- et vaut-il mieux freiner pour en payer moins ?}
\concret
\begin{pts}
\item \textbf{La vitesse} d'une méthode se mesure à chaque coupe :
$\ \mathrm{TO}_t=\tfrac12\sum_i\bigl|w^{\text{réalisé}}_i-w^{\text{dérivé}}_i\bigr|$, la fraction du portefeuille
qui change de mains. \textbf{Ce n'est pas un réglage, c'est une conséquence de la fenêtre} : sur 3 ans le score
bouge à peine d'un mois à l'autre, sur 42 jours il se reconstruit.
\item \textbf{Les frais} : on ampute la valeur à chaque coupe, $\ V_t\leftarrow V_t(1-2\,\mathrm{TO}_t\,c)$ avec
$c=10$~pb. Le \textbf{facteur 2} parce que $\mathrm{TO}_t$ est le montant acheté, égal au montant vendu : on paie
des deux côtés.
\end{pts}
\begin{center}\scriptsize
\begin{tabular}{@{}l r r r r@{}}
\toprule
 & rotation/an & coût/an & excès net & $\alpha_4$ net \\
\midrule
M1 \emph{(commune)} & \textbf{715~\%} & $1{,}59$~pt & $\mathbf{-1{,}24}$~\% & $-0{,}46$ \\
M2 --- NANC & 63~\% & $0{,}18$ & $+1{,}17$ & $+0{,}80$ \\
\textbf{M3 --- NANC} & 65~\% & $0{,}18$ & $+2{,}80$ & $+1{,}59$ \\
M4 --- NANC & 64~\% & $0{,}19$ & $+4{,}37$ & $+1{,}58$ \\
M2 --- GOP & 79~\% & $0{,}19$ & $+0{,}32$ & $+1{,}11$ \\
\textbf{M3 --- GOP} & 68~\% & $0{,}19$ & $+1{,}04$ & $+1{,}51$ \\
M4 --- GOP & 71~\% & $0{,}20$ & $+1{,}32$ & $+2{,}10$ \\
\bottomrule
\end{tabular}
\end{center}
\obs{\textbf{Première réponse : le frottement décide du sort de M1.} Onze fois plus rapide que les autres, elle
perd $1{,}59$~pt/an et \emph{passe sous zéro} --- son excès n'aurait pas survécu au passage d'ordre. Les
pondérations au dollar, elles, ne laissent qu'un cinquième de point.}
\obs{\textbf{Reste le frein} de l'étape 6 : au lieu de rejoindre la cible, on n'en parcourt qu'une fraction,
$\ \lambda_R=\min(1,\tau_{\max}/\mathrm{TO}_t)$, puis $w^{\text{réalisé}}=w^{\text{dérivé}}+\lambda_R
(w^{\text{cible}}-w^{\text{dérivé}})$. À $\tau_{\max}=25\,\%$ il mord \textbf{à chaque coupe} sur M1 --- mais sur
\textbf{3~\% des coupes} seulement sur M3. Pour que la question ait un sens, il faut descendre bien plus bas.}
\fig{meth_frein.png}{le frein appliqué à M3, \textbf{frais compris}. À 2~\%/mois la contrainte sature : la rotation
tombe à 24~\%/an des deux côtés et le portefeuille ne suit plus sa cible --- ce n'est plus M3 ralentie, c'est une
inertie qui laisse courir les positions.}
\lecon{seconde réponse --- \textbf{le frein n'est pas retenu}. Il gagne $+0{,}71$~pt sur NANC et perd $-0{,}85$~pt
sur GOP. Un réglage qui aide un camp et détruit l'autre est exactement ce que la règle de décision écarte.}

% ═══════════ VERDICT ═══════════
\meth{$=$}{La méthode retenue, et pourquoi c'est M3}{on choisit \emph{avant} de regarder les rendements : la plus
parcimonieuse dont le résultat tient sur les \textbf{deux} partis.}
\begin{pts}
\item \textbf{M3}, pas M4 : M4 gagne plus mais son $t$ est plus faible sur les deux camps, son IC double, et il
coûte un paramètre de plus. \textbf{On ne retient pas la méthode qui a le mieux marché, on retient la plus sobre
qui marche des deux côtés.}
\item M3 est la seule dont le $t$ dépasse 1,4 \textbf{et} pour NANC \textbf{et} pour GOP --- $1{,}47$ et $1{,}62$
en excès annuel, $1{,}81$ et $1{,}82$ en $\alpha_4$. La cohérence entre deux échantillons \emph{distincts} vaut
plus qu'un chiffre isolé --- sans être deux preuves indépendantes : les deux paniers partagent des titres et la
même période.
\end{pts}
\[ \boxed{\begin{gathered}
s_i=B_i\,m_i\ \text{ sur 3 ans à }\delta,\quad \mathcal U=\text{top-150}\\[0.4mm]
w_i=\min(c,\lambda v_i),\quad c=8{,}5\,\%
\end{gathered}} \]
\fig{meth_nav.png}{les NAV, \textbf{toutes rebasées à 100 au 03/02/2014} --- la seule façon de les comparer d'un seul regard (échelle log). Sur cette fenêtre commune : M1 \textbf{530}, M2 \textbf{584}, \textbf{M3 672}, M4 \textbf{698}, \textbf{SPY 528}.}

% ═══════════ CE QUI MANQUE ═══════════
{\color{bleu}\bfseries $\varnothing$ $\cdot$ Ce qu'aucune de ces méthodes ne fait}\par\vspace{0.3mm}
{\small \textbf{Vendre à découvert} --- testé, catastrophique : le panier des titres vendus \emph{est} le marché
($\beta=-1{,}03$), et une jambe courte seule termine à 16 sur base 100 --- et cela \emph{à $\tau$}, la date la
plus favorable. \textbf{Suivre les options, les ETF ou les
obligations} --- seules les actions entrent. \textbf{Payer l'impact de marché} --- le passage d'ordre est facturé
(bloc \textdollar), l'effet d'un ordre \emph{sur le prix} non. \textbf{Sortir parce qu'un élu vend} --- sauf M1,
et seulement sous deux ans.\par}

{\color{bleu}\bfseries $\neq$ $\cdot$ Ce qui empêche de conclure, et comment le lever}\par\vspace{0.3mm}
\begin{pts}
\item \textbf{La significativité.} Aux seuils usuels, l'$\alpha_4$ de M3 n'est \textbf{pas} significatif
($t=1{,}81$ et $1{,}82$ contre $1{,}96$) ; celui de M4 côté GOP l'est \textbf{marginalement} ($1{,}99$).
\textbf{Ce qui manque est une validation hors échantillon} : M3 n'a jamais été testée sur une période réservée,
et c'est l'étape qui transformerait un résultat encourageant en résultat établi.
\item \textbf{La capacité.} Le passage d'ordre est facturé (bloc \textdollar) et ne coûte que $0{,}18$~pt/an
à M3. Reste l'\emph{impact} : un fonds assez gros déplacerait le prix des lignes qu'il achète.
\item \textbf{La période.} 2014--2026 est un marché haussier dominé par la croissance, précisément le tilt de M3.
\item \textbf{Où loge l'$\alpha$ ?} \emph{Dans le temps}, c'est fait (bloc $\circlearrowright$) : diffus, positif
$\approx$70~\% des fenêtres. \emph{En coupe} --- par secteur, par taille de dépôt, par élu --- ce ne l'est pas.
Qu'il se concentre ou qu'il reste diffus, les deux réponses valent d'être sues.
\end{pts}
\vspace{1mm}\hrule\vspace{1mm}
{\scriptsize\color{grisf} Tous les chiffres sont des sorties du notebook \path{11_Strategie_Ticker.ipynb} :
\S16 (les quatre méthodes, runs déclarés avant exécution), \S16.1 (les quatre facteurs, validés sur le SPY :
$\beta_{\mathrm M}=0{,}980$, $\alpha=+0{,}04$~\%/an), \S16.3 ($\alpha$ et $\beta$ glissants), \S16.4 (rotation,
frais, frein), \S16.5 (audit de méthode). Facteurs Fama-French-Carhart, \path{cache/ff_factors.csv}.
\textbf{Mesure indépendante} : Baulkaran \& Jain, \emph{Economics Letters} 250 (2025), régressent les
\emph{ETF réels} NANC et KRUZ sur \textbf{cinq} facteurs (les quatre mêmes plus CMA), en rendement brut et sur
2023--24 seulement, et ne trouvent aucun $\alpha$ significatif. Objet différent du nôtre --- le fonds
commercialisé, non une stratégie reconstruite --- les deux résultats se lisent en parallèle. L'étude de fidélité
à ce fonds, dont M4 reprend le filtre, fait l'objet d'un document distinct, \path{FICHE_NANC_GOP.pdf}.\par}
\end{multicols}



\end{document}
